% =========================================================================
% Metaheuristics Project (31-MM34, SoSe 26) -- Example_V0
% Compile: pdflatex Example_V0.tex
% =========================================================================
\documentclass[10pt, a4paper]{article}

\usepackage[utf8]{inputenc}
\usepackage[T1]{fontenc}
\usepackage{lmodern}
\usepackage[margin=2.5cm]{geometry}
\usepackage{amsmath,amssymb}
\usepackage{booktabs}
\usepackage[hidelinks]{hyperref}

\newcommand{\ub}{\%\,UB}

\title{An Adaptive Large Neighborhood Search for the\\
       Prize-Collecting Skill Vehicle Routing Problem}
\author{Umais Chaudhary, Rehman Kerimov, Dustin Schulz, Marco Leander vom Orde \\ \small 31-MM34 Metaheuristics, SoSe 26, Bielefeld University}
\date{Juli 2026}

\begin{document}
\maketitle

\begin{abstract}
We present a metaheuristic for the prize-collecting Skill Vehicle Routing
Problem, in which a small team of couriers with heterogeneous skill sets
serves an optional subset of customers under time windows and shift limits so
as to maximize collected profit. Guided by the literature on team orienteering
and technician routing, we adopt an Adaptive Large Neighborhood Search (ALNS)
with simulated-annealing acceptance. Our central observation is that, under
the project's 30-second CPython budget, \emph{move-evaluation throughput} ---
not operator sophistication --- is the binding constraint: constant-time
insertion feasibility via forward time slacks and a per-route move cache
raise the iteration count on the largest instance by a factor of 250 and turn
an otherwise idle framework into a competitive solver. On the 20 benchmark
instances (50--1000 customers) the final configuration attains 63.0\,\% of a
singleton-feasibility upper bound as a mean over three seeds, with all 60
runs feasible and within the time limit, matching the best single run of a
heavily tuned single-trajectory LNS baseline.
\end{abstract}

% =========================================================================
\section{Introduction}

A pharmacy plans daily home-delivery tours for a small team of couriers.
Deliveries differ in handling requirements (refrigeration, controlled
substances, device setup), and each courier is certified for only a subset of
these \emph{skills}. Because the team is small, not every request can be
served: each served request yields a profit, and the objective is to select
and route a subset of customers maximizing total profit while respecting
skill compatibility, customer time windows and courier shifts.

Three properties shape the algorithmic landscape. First, the problem is
\emph{selective}: customer choice and routing must be optimized jointly, which
weakens methods that rely on all customers being present. Second, skills act
as hard compatibility filters between customers and couriers, so courier
capacity is a heterogeneous, contested resource. Third, the project imposes a
hard wall-clock budget (e.g., 30\,s) on pure-Python code, so the effective
search effort is limited by how cheaply a candidate move can be evaluated.

Our design follows a single guiding principle that also structures this
report: \emph{evidence first, efficiency second, operator richness third}.
The literature on the two closest problem families --- the team orienteering
problem with time windows (TOPTW) and technician routing --- consistently
favors destroy-and-repair searches over population methods
(Section~\ref{sec:related}); we therefore adopt ALNS
\cite{ropke2006,pisinger2007}. A first, faithful implementation of that
framework, however, completed only \emph{four} iterations within 30\,s on the
largest instance --- less than the construction heuristic alone achieves.
Section~\ref{sec:method} describes how constant-time feasibility checking
\cite{savelsbergh1992}, a per-route move cache and a wall-clock-based
annealing schedule convert the same framework into one that performs roughly
1{,}000 iterations on 1000 customers and 50{,}000 on 50 customers.
Sections~\ref{sec:components} and~\ref{sec:overall} experimentally justify
each component and evaluate the final solver.

Our contributions are: (i) an ALNS for the prize-collecting Skill VRP with a
skill-aware operator portfolio, including a problem-specific
\emph{skill-scarcity} destroy operator; (ii) an evaluation architecture that
makes the framework viable in interpreted Python (forward time slacks, move
caching, early-exit delta evaluation), validated by differential testing;
(iii) a systematic component study showing, among other things, that destroy
\emph{size} is best treated as part of the operator identity learned by the
adaptive weights; and (iv) a noise-robust final evaluation (three seeds per
instance) reaching 63.0\,\ub{} on average with zero infeasible or late runs.

% =========================================================================
\section{Problem Statement}\label{sec:problem}

Let $C=\{1,\dots,n\}$ be the customers and $B$ the couriers; node $0$ is the
pharmacy. Each customer $c$ has profit $p_c$, service time $s_c$, time window
$[e_c,\ell_c]$ and required skills $R_c\subseteq\{1,\dots,S\}$. Each courier
$b$ has a shift $[\mathit{start}_b,\mathit{end}_b]$ and skills
$Q_b\subseteq\{1,\dots,S\}$. Travel times $d_{ij}$ are Euclidean distances
rounded to the nearest integer, identical for all couriers. A plan assigns to
every courier a tour starting and ending at the pharmacy within its shift; a
courier may serve $c$ only if $R_c\subseteq Q_b$; arrival at $c$ must not
exceed $\ell_c$ (early arrivals wait until $e_c$); each customer is visited at
most once. With $x_c\in\{0,1\}$ indicating service, the objective is
\begin{equation}
  \max \sum_{c\in C} p_c\, x_c .
  \label{eq:objective}
\end{equation}
The problem contains the team orienteering problem and hence the TSP as
special cases and is NP-hard. Since optima are unknown, we report solution
quality relative to the \emph{singleton upper bound}
$\mathrm{UB}=\sum_{c\,:\,\exists b:\, c \text{ servable solo by } b} p_c$,
the profit of all customers that at least one courier could serve as a
one-customer tour. UB ignores all interaction between customers and is
therefore loose; it is nevertheless instance-independent of the solver and
suitable for relative comparisons.

% =========================================================================
\section{Related Work}\label{sec:related}

\subsection{Problem classification}
The selective aspect --- optional customers, profit maximization under
routing-time restrictions --- places the problem in the orienteering family
\cite{vansteenwegen2011}; with multiple vehicles and hard time windows its
routing core is the TOPTW \cite{labadie2012,gunawan2016}. The skill aspect
corresponds to the Skill VRP of Cappanera et al.~\cite{cappanera2011} and the
technician routing and scheduling problem (TRSP)
\cite{kovacs2012,pillac2013}, which combines skills, time windows and worker
shifts exactly as here.

\subsection{Metaheuristics for selective routing with time windows}
The TOPTW literature is dominated by trajectory-based search: the fast
iterated local search of Vansteenwegen et al.~\cite{vansteenwegen2009}
remains the low-runtime reference; granular VNS \cite{labadie2012}, the
three-component I3CH \cite{hu2014} and a well-tuned ILS/SA hybrid
\cite{gunawan2017} define the quality frontier. A shared efficiency idea is
constant-time insertion feasibility via forward time slacks
\cite{savelsbergh1992,vansteenwegen2009}, which we adopt.

\subsection{LNS, ALNS and population methods}
Large neighborhood search \cite{shaw1998} and its adaptive variant
\cite{ropke2006,pisinger2007} repeatedly destroy part of a solution (random,
related, worst removal) and repair it by greedy or regret insertion; SISR
\cite{christiaens2020} shows that carefully designed removals keep this
family state of the art. For the TRSP --- the closest published relative of
our problem --- Kovacs et al.~\cite{kovacs2012} report ALNS clearly
outperforming exact approaches beyond small instances; the home-health-care
survey of Ciss\'e et al.~\cite{cisse2017} reaches the same conclusion for
compatibility-constrained routing with optional visits. Hybrid genetic search
\cite{vidal2012,vidal2022} defines the state of the art for classical VRPs,
but its strengths (giant-tour decoding, dense route-improvement
neighborhoods) degrade when customer \emph{selection} is part of the problem
and hard time windows, shifts and skill filters render most offspring
infeasible.

\subsection{Synthesis}
Both closest problem families point to the same design: destroy-and-repair
search with insertion-based repair, feasibility filtering for skills, and an
annealing-type acceptance --- i.e., ALNS. Under a short, interpreted-Python
time budget the cheap-move efficiency of ALNS weighs even more heavily
against population methods. This is the method we implement; the remainder of
the paper shows that realizing its potential in Python is chiefly an
\emph{evaluation-efficiency} problem.

% =========================================================================
\section{Solution Approach}\label{sec:method}

\subsection{Framework overview}
Algorithm~1 shows the search. After preprocessing and a greedy
multi-start construction, the ALNS iterates destroy--repair cycles; operators
are drawn with probability proportional to adaptive weights; candidates are
accepted by simulated annealing whose temperature decays with \emph{elapsed
wall-clock time}; new global bests trigger an Or-opt polish that compresses
routes and refills freed capacity. The best feasible solution is printed; an
empty plan is the always-feasible fallback.

\begin{figure}[t]
\hrule\vspace{3pt}
\noindent\textbf{Algorithm 1:} ALNS for the prize-collecting Skill VRP
\vspace{3pt}\hrule\vspace{4pt}
\small
\begin{tabbing}
\hspace{1.5em}\=\hspace{1.5em}\=\hspace{1.5em}\=\kill
1:\> $x \gets$ best of four greedy insertion orderings;\quad $x^{*} \gets x$\\
2:\> \textbf{while} $\mathrm{now} < \mathrm{deadline}$ \textbf{do}\\
3:\>\> draw destroy $d$ and repair $r$ proportional to weights;\quad
        $x' \gets r(d(x))$ \hfill$\triangleright$ feasible insertions only\\
4:\>\> $T \gets T_0\,(T_{\min}/T_0)^{(\mathrm{now}-t_0)/(\mathrm{deadline}-t_0)}$
        \hfill$\triangleright$ time-based cooling\\
5:\>\> \textbf{if} $f(x') \ge f(x)$ \textbf{or}
        $\mathrm{rand} < e^{(f(x')-f(x))/T}$ \textbf{then} $x \gets x'$\\
6:\>\> \textbf{if} $f(x') > f(x^{*})$ \textbf{then}
        $x^{*} \gets \textsc{PolishAndFill}(x')$;\quad $x \gets x^{*}$
        \hfill$\triangleright$ Or-opt + refill, rate-limited\\
7:\>\> update operator scores;\quad every 100 iterations update weights\\
8:\>\> return to $x^{*}$ after 2000 non-improving iterations;\quad
        restart and reheat after 500 rejections\\
9:\> \textbf{return} $x^{*}$
\end{tabbing}
\vspace{-6pt}\hrule
\end{figure}

\subsection{Preprocessing and construction heuristic}
Preprocessing computes the rounded distance matrix, for every customer the
list of couriers that can serve it \emph{solo} (skill subset test plus a
singleton-tour time check) --- customers with an empty list are discarded once
and for all --- and a profit-density lower bound $p_c/\tau_c$, where $\tau_c$
estimates the minimal time to serve $c$. Construction inserts customers one
at a time at the feasible position of minimal travel detour over all
compatible couriers. Four orderings (profit; profit density; fewest feasible
couriers first; a lexicographic hybrid) are all evaluated and the most
profitable start is kept; on the largest instance this multi-start costs
under 5\,s and the gap between orderings reaches 20\,\% (Table~\ref{tab:final}
shows the selected starts).

\subsection{Destroy operators}
All destroy sizes $q$ are drawn per call from a fraction range of currently
served customers. The portfolio combines general-purpose and
problem-specific operators:
\emph{random removal} (diversification);
\emph{related removal} \cite{shaw1998}: a random seed customer and its
nearest served neighbors, removing a coherent region a repair can re-route;
\emph{worst-detour removal} \cite{ropke2006}: customers with the lowest
profit per travel detour in their current position, biased sampling;
\emph{worst-density removal}: the static analogue based on $p_c/\tau_c$;
and \emph{skill-scarcity removal}, which targets low-profit customers
occupying couriers whose skill set is in high demand among the currently
unserved customers --- freeing contested skill capacity that geometric
operators cannot ``see''.
Because no single destroy size suits all instances
(Section~\ref{sec:exp-size}), the three geometric operators exist in a
\emph{small} ($U(0.03,0.12)$, cap 40) and a \emph{large} ($U(0.12,0.28)$, cap
70) variant, making the size part of the operator identity that the adaptive
weights learn per instance.

\subsection{Repair operators}
Repair candidates are the removed customers plus 100 unserved extras --- half
the top profit-density customers, half sampled uniformly, since a purely
static candidate list would starve the remaining pool. All repairs generate
\emph{feasible} insertions only.
\emph{Greedy best-first insertion} repeatedly applies the globally best
insertion (highest profit, ties by minimal detour);
\emph{regret-2 insertion} \cite{ropke2006} prioritizes customers whose
second-best slot is much worse than their best;
\emph{sequential cheapest insertion} \cite{kovacs2012} performs a single pass
over the candidates (profit or random order), inserting each immediately at
its cheapest feasible position --- roughly three times cheaper per call, it
trades per-repair care for iteration throughput;
a \emph{noised} sequential variant perturbs the position ranking by
$\pm 0.1\,d_{\max}$ \cite{ropke2006} to escape deterministic ruts.

\subsection{Efficient move evaluation}\label{sec:eval}
This subsection is the technical core. Three mechanisms reduce the cost of
the dominant operation --- ``what is the best feasible insertion of customer
$c$ into route $r$?'' --- by orders of magnitude.

\textbf{(i) Forward time slacks.} Every route caches arrival, service-start
and departure times plus the forward time slack \cite{savelsbergh1992}:
$\mathit{slack}_i$ is the maximal delay of the service start at position $i$
that keeps positions $i,\dots,m$ and the depot return feasible, computed
backwards in $O(m)$ when the route changes. An insertion at position $i$ is
then feasible iff the arrival at the new customer meets its due time and the
induced delay of the successor does not exceed $\mathit{slack}_i$ --- an
$O(1)$ test replacing an $O(m)$ re-simulation. Two monotonicity pruning rules
skip positions that cannot host the customer at all.

\textbf{(ii) Move caching.} Inserting a customer modifies exactly one route,
so cached best moves of all candidates on \emph{other} routes remain exact.
The repair operators therefore build a per-(customer, vehicle) cache of top
moves once and, after each insertion, recompute only the entries of the
modified vehicle \cite{ropke2006}. Previously each single insertion
re-evaluated every candidate on every vehicle.

\textbf{(iii) Early-exit delta evaluation.} For the rare evaluation of
infeasible routes (the penalty branch of the SA), the route suffix is
re-simulated only until its timeline re-synchronizes with the cached one;
the cached suffix penalty then applies verbatim.

Correctness of both evaluation paths was established by differential testing
against brute-force route re-simulation ($\approx$16{,}000 random cases, zero
deviations). Section~\ref{sec:exp-throughput} quantifies the effect: a factor
of $\approx$250 in iteration throughput on the largest instance.

\subsection{Acceptance, adaptivity and intensification}
Two calibration insights matter more than exact parameter values. First,
the initial temperature must be scaled to the magnitude of a \emph{move}:
destroy--repair deltas are sums of a few customer profits (10--100)
regardless of instance size, whereas total profit grows twentyfold from
$n{=}50$ to $n{=}1000$; scaling $T_0$ to the solution value (as often done)
turns the search into a random walk on large instances. We use
$T_0 = 1.5\times$ mean customer profit. Second, cooling must be
\emph{time-based}: iteration throughput spans two orders of magnitude across
the benchmark, so a fixed geometric rate per iteration is either frozen on
small or never cools on large instances; we decay $T$ exponentially in the
elapsed fraction of the budget (Algorithm~1, line~4).
Operator weights are updated every 100 iterations from scores (new global
best 25, improving 10, accepted 3, rejected 0), smoothed with reaction factor
0.2 and normalized by the operator's measured runtime, so expensive operators
must earn their cost. The current solution returns to the incumbent best
after 2000 non-improving iterations; after 500 consecutive rejections the
search restarts from a randomized greedy solution and reheats.

\subsection{Local-search polish}
Insertion-based repairs never reorder a route, leaving travel time on the
table. Whenever a new global best appears (rate-limited to twice per second),
every route is compressed by \emph{Or-opt}: segments of one to three
consecutive customers are relocated within the route if this strictly reduces
travel and keeps the route feasible (first improvement, repeated to a fixed
point). Or-opt changes no profit by itself; its value is the freed shift
capacity, which a subsequent greedy pass immediately refills with additional
customers. The search then continues from the compressed solution.

% =========================================================================
\section{Experimental Investigation of the Components}\label{sec:components}

\subsection{Setup and evaluation metric}
All experiments: one core of an Apple-silicon laptop, CPython~3.9, no
third-party packages, 30\,s budget (component studies use 25\,s), seed~1
unless stated. Quality is reported as profit or as \ub{}, the percentage of
the singleton upper bound of Section~\ref{sec:problem}. Every reported
solution passes the provided checker. Because single runs scatter by roughly
$\pm$1\,\% per instance, we treat sub-percent single-seed differences as
noise and base final conclusions on three seeds
(Section~\ref{sec:overall}).

\subsection{Move-evaluation throughput}\label{sec:exp-throughput}
Table~\ref{tab:throughput} isolates the evaluation mechanisms of
Section~\ref{sec:eval} on the largest instance. The naive framework cannot
even leave its construction start: four iterations consume the entire
budget. Quality tracks throughput almost mechanically --- the search
components only become \emph{observable} once evaluation is cheap. This
justifies our efficiency-first design order.

\begin{table}[t]
\centering\small
\caption{Throughput ablation, instance n1000 (30\,s, seed 1).}
\label{tab:throughput}
\begin{tabular}{lrr}
\toprule
evaluation scheme & iterations & best profit \\
\midrule
naive suffix re-simulation                   & 4        & 30\,193 \\
\; + early-exit deltas + move cache          & 517      & 33\,526 \\
\; + slack fast path                          & $\approx$1\,000 & 35\,000+ \\
\; + full final portfolio and polish          & 962      & 36\,239 \\
\bottomrule
\end{tabular}
\end{table}

\subsection{Destroy size}\label{sec:exp-size}
Fixed destroy fractions are not robust: 15\,\% wins on n1000 but loses on
n400, 6\,\% vice versa (Table~\ref{tab:size}, left). A single randomized
range helps but still trades sizes off between instances. Making the size
part of the \emph{operator identity} --- small and large variants of each
geometric destroy, selected by the adaptive weights --- dominates both
alternatives across seeds (right). The learned usage shares confirm the
mechanism: on n400 the small variants are drawn 2--3$\times$ more often than
the large ones, on n1000 the preference narrows.

\begin{table}[t]
\centering\small
\caption{Left: fixed vs.\ randomized destroy fraction (25\,s). Right:
size-split operator variants, two seeds (25\,s).}
\label{tab:size}
\begin{tabular}{lrrr@{\qquad}rr}
\toprule
 & \multicolumn{3}{c}{single range} & \multicolumn{2}{c}{size-split (s1 / s2)} \\
instance & 0.15 & 0.06 & $U(0.04,0.18)$ & \multicolumn{2}{c}{} \\
\midrule
n100  & --      & --      & 2\,910\,$^{a}$ & \multicolumn{2}{c}{3\,020 / 3\,011} \\
n400  & 12\,475 & 12\,824 & 12\,913 & \multicolumn{2}{c}{--} \\
n600  & --      & --      & 20\,750$^{a}$ & \multicolumn{2}{c}{21\,098 / 20\,915} \\
n1000 & 35\,355 & 34\,838 & 35\,032 & \multicolumn{2}{c}{36\,216 / 35\,929} \\
\bottomrule
\multicolumn{6}{l}{\footnotesize $^{a}$wide range $U(0.05,0.28)$, 30\,s benchmark configuration.}
\end{tabular}
\end{table}

\subsection{Operator portfolio}
\textbf{Correlated and solution-dependent removal.} The initial portfolio
(random, static worst-density, skill-scarcity) lacked spatially correlated
removals; the static density ranking moreover removes the same customers
every iteration. Adding related removal and worst-detour removal improved
every instance tested (n100: 2\,823$\to$2\,876; n250: 9\,450$\to$9\,787;
n800: 24\,829$\to$25\,713) --- on mid and large instances the repair can only
restructure a region if the destroy frees a coherent one.

\textbf{Cheap sequential repair.} The best-first repairs re-select the global
best candidate after every insertion; adding the single-pass sequential
insertion roughly triples the attainable iteration count and let the adaptive
weights trade care against throughput (n250: 9\,694$\to$9\,976; n1000:
35\,610$\to$35\,808). Its usage share ends highest on \emph{every} instance
(40--46\,\%), while regret-2 keeps a minor share --- adaptivity working as
intended.

\textbf{Candidate pool.} Repairing from the full unserved pool instead of the
sampled 100 extras was tested and \emph{rejected}: on n850 the iteration
count collapses (759$\to$214) and about 1\,000 profit is lost. At this budget
throughput beats per-repair completeness.

\subsection{Acceptance design}
Replacing the solution-scaled $T_0$ by the move-scaled one and the
iteration-based by the time-based schedule was worth roughly $+1.5$ points
\ub{} on large instances (the n1000 best rose from 33\,526 to
$\approx$35\,000 with identical operators); subsequent factor sweeps
($T_0\in\{1,1.5,2\}\times$ mean profit) changed results by only about
$\pm$1\,\%, i.e., the \emph{scale} matters, not the exact factor. The
return-to-best intensification gave small but consistently non-negative
changes.

\subsection{Local-search polish and noised insertion}
The final two components added $+0.3$ points \ub{} on average (single-seed
spot checks: n100 2\,946$\to$3\,011, n1000 36\,016$\to$36\,239). Polishing
pays because freed travel time converts into additional served customers via
the refill pass; noise removes deterministic repair ruts at negligible cost.

% =========================================================================
\section{Overall Performance}\label{sec:overall}

\subsection{Final results}
Table~\ref{tab:final} reports the final configuration on all 20 benchmark
instances with \textbf{three seeds each} (60 runs): every run passed the
checker and finished within the limit (max.\ wall time 29.2\,s). The mean
over seed means is \textbf{63.0\,\ub}. Seed spread is small (typically below
1.5\,\% between min and max), so this figure is an \emph{expected} quality,
not a lucky draw. The ALNS improves its greedy start by 17--36\,\%,
underlining that the search --- not the construction --- carries the result.

\begin{table}[t]
\centering\small
\caption{Final results: 30\,s, three seeds per instance. Greedy = multi-start
construction; ALNS = mean profit over seeds; LNS = single-run baseline of
Section~\ref{sec:baseline}.}
\label{tab:final}
\begin{tabular}{lrrrrrrr}
\toprule
instance & $n$ & $B$ & UB & greedy & ALNS & \ub & LNS \\
\midrule
n50\_s3\_k0.0   &   50 &  2 &  2\,717 &  1\,501 &  1\,796 & 66.1 &  1\,793 \\
n75\_s4\_k0.5   &   75 &  2 &  4\,255 &  2\,058 &  2\,790 & 65.6 &  2\,784 \\
n100\_s5\_k1.0  &  100 &  2 &  5\,160 &  2\,292 &  2\,996 & 58.1 &  3\,004 \\
n125\_s6\_k1.5  &  125 &  3 &  6\,866 &  3\,697 &  4\,318 & 62.9 &  4\,223 \\
n150\_s8\_k2.0  &  150 &  3 &  8\,259 &  4\,065 &  4\,861 & 58.9 &  4\,936 \\
n200\_s3\_k0.0  &  200 &  5 & 10\,236 &  5\,504 &  7\,454 & 72.8 &  7\,487 \\
n250\_s4\_k0.5  &  250 &  6 & 14\,034 &  7\,845 &  9\,795 & 69.8 &  9\,874 \\
n300\_s5\_k1.0  &  300 &  7 & 16\,729 &  8\,331 & 11\,248 & 67.2 & 11\,450 \\
n350\_s6\_k1.5  &  350 &  8 & 19\,050 &  9\,708 & 12\,581 & 66.0 & 12\,583 \\
n400\_s8\_k2.0  &  400 & 10 & 22\,335 & 10\,426 & 12\,942 & 57.9 & 12\,741 \\
n450\_s3\_k0.0  &  450 & 11 & 24\,478 & 11\,003 & 14\,271 & 58.3 & 14\,155 \\
n500\_s4\_k0.5  &  500 & 12 & 27\,820 & 13\,344 & 17\,690 & 63.6 & 17\,620 \\
n600\_s5\_k1.0  &  600 & 15 & 32\,477 & 16\,277 & 20\,734 & 63.8 & 21\,177 \\
n700\_s6\_k1.5  &  700 & 17 & 37\,852 & 18\,732 & 23\,779 & 62.8 & 23\,758 \\
n750\_s8\_k2.0  &  750 & 18 & 40\,329 & 19\,542 & 25\,376 & 62.9 & 25\,217 \\
n800\_s3\_k0.0  &  800 & 20 & 43\,128 & 20\,995 & 25\,806 & 59.8 & 25\,491 \\
n850\_s4\_k0.5  &  850 & 21 & 46\,190 & 21\,708 & 29\,629 & 64.1 & 29\,875 \\
n900\_s5\_k1.0  &  900 & 22 & 49\,703 & 23\,101 & 28\,714 & 57.8 & 28\,660 \\
n950\_s6\_k1.5  &  950 & 23 & 52\,946 & 22\,960 & 29\,391 & 55.5 & 29\,349 \\
n1000\_s8\_k2.0 & 1000 & 25 & 54\,664 & 30\,193 & 36\,211 & 66.2 & 35\,781 \\
\midrule
average &&&&&& \textbf{63.0} & \\
\bottomrule
\end{tabular}
\end{table}

\subsection{Comparison with a tuned single-trajectory LNS}\label{sec:baseline}
As an internal yardstick we compare against an independently implemented,
heavily optimized plain LNS developed earlier in the project (greedy
construction; random, worst and related removal; full-pool greedy insertion;
record-to-record-travel acceptance). Its best benchmark log reaches
63.0\,\ub{} --- as a \emph{single run}. Our three-seed mean matches this
figure, i.e., the expected quality of the ALNS equals the best observed run
of the baseline. Head-to-head (our seed mean vs.\ the baseline's single run)
the ALNS wins or ties 12 of 20 instances; on the largest instance every
individual seed beats the baseline (36\,111--36\,268 vs.\ 35\,781). The
comparison also validates the framework overhead: eight destroy and five
repair operators with adaptive selection are fully amortized by the $O(1)$
evaluation.

\subsection{Development progression and scaling}
Table~\ref{tab:progress} summarizes the average quality after each
development stage --- a compact confirmation of the report's red thread:
throughput first (61.0), then operator richness (62.2--62.7), then local
search (63.0). Iteration counts at 30\,s range from $\approx$50{,}000
($n{=}50$) to $\approx$1{,}000 ($n{=}1000$), yet \ub{} shows no systematic
degradation with $n$: the hardest instances are those with scarce skills
(suffixes k1.5/k2.0) and few couriers per skill, not the largest ones ---
consistent with skill capacity, not geometry, being the binding resource.
The solver degrades gracefully under small budgets: at 2\,s on $n{=}1000$ it
still returns a feasible (partial greedy) solution within the limit.

\begin{table}[t]
\centering\small
\caption{Average \ub{} over all 20 instances after each development stage
(30\,s; final row: mean over three seeds).}
\label{tab:progress}
\begin{tabular}{lr}
\toprule
stage & avg.\ \ub \\
\midrule
faithful ALNS framework (naive evaluation)              & n/a (idle on large $n$) \\
+ slack fast path, move cache, time-based SA            & 61.0 \\
+ related \& worst-detour removal, candidate sampling   & 62.2 \\
+ sequential repair, size-split destroy operators       & 62.7 \\
+ Or-opt polish \& refill, noised insertion (final)     & \textbf{63.0} \\
\bottomrule
\end{tabular}
\end{table}

% =========================================================================
\section{Conclusion}

We solved the prize-collecting Skill VRP with an ALNS whose design was driven
by one empirical lesson: under a short wall-clock budget in interpreted
Python, evaluation throughput is the primary bottleneck, and every search
refinement only becomes effective after it is removed. Constant-time
insertion feasibility via forward time slacks, per-route move caching and a
wall-clock-based annealing schedule raised throughput by two orders of
magnitude; on this foundation, a skill-aware operator portfolio with
size-split destroy variants, cheap sequential repairs, and an Or-opt polish
with capacity refill lifted the average quality to 63.0\,\% of the singleton
upper bound as a three-seed mean --- with all 60 final runs feasible and on
time, and matching the best single run of a heavily tuned LNS baseline.
Promising next steps are string-based removals in the spirit of SISR
\cite{christiaens2020}, a time-aware insertion score $p_c - \lambda\cdot
\Delta\mathrm{travel}$ reflecting that shift time is the scarce resource, and
a set-covering post-optimization over the route pool collected during the
search \cite{hu2014}.

% =========================================================================
\begin{thebibliography}{19}\small

\bibitem{cappanera2011} P.~Cappanera, L.~Gouveia, M.G.~Scutell\`a.
The skill vehicle routing problem. \emph{Network Optimization} (INOC 2011),
LNCS 6701, 354--364, 2011.

\bibitem{christiaens2020} J.~Christiaens, G.~Vanden Berghe.
Slack induction by string removals for vehicle routing problems.
\emph{Transportation Science} 54(2), 417--433, 2020.

\bibitem{cisse2017} M.~Ciss\'e, S.~Yal\c{c}{\i}nda\u{g}, Y.~Kergosien,
E.~\c{S}ahin, C.~Lent\'e, A.~Matta. OR problems related to home health care:
a review of relevant routing and scheduling problems.
\emph{Operations Research for Health Care} 13--14, 1--22, 2017.

\bibitem{gunawan2016} A.~Gunawan, H.C.~Lau, P.~Vansteenwegen.
Orienteering problem: a survey of recent variants, solution approaches and
applications. \emph{European Journal of Operational Research} 255(2),
315--332, 2016.

\bibitem{gunawan2017} A.~Gunawan, H.C.~Lau, P.~Vansteenwegen, K.~Lu.
Well-tuned algorithms for the team orienteering problem with time windows.
\emph{Journal of the Operational Research Society} 68(8), 861--876, 2017.

\bibitem{hu2014} Q.~Hu, A.~Lim. An iterative three-component heuristic for
the team orienteering problem with time windows.
\emph{European Journal of Operational Research} 232(2), 276--286, 2014.

\bibitem{kovacs2012} A.A.~Kovacs, S.N.~Parragh, K.F.~Doerner, R.F.~Hartl.
Adaptive large neighborhood search for service technician routing and
scheduling problems. \emph{Journal of Scheduling} 15(5), 579--600, 2012.

\bibitem{labadie2012} N.~Labadie, R.~Mansini, J.~Melechovsk\'y,
R.~Wolfler Calvo. The team orienteering problem with time windows: an
LP-based granular variable neighborhood search.
\emph{European Journal of Operational Research} 220(1), 15--27, 2012.

\bibitem{pillac2013} V.~Pillac, C.~Gu\'eret, A.L.~Medaglia. A parallel
matheuristic for the technician routing and scheduling problem.
\emph{Optimization Letters} 7(7), 1525--1535, 2013.

\bibitem{pisinger2007} D.~Pisinger, S.~Ropke. A general heuristic for
vehicle routing problems. \emph{Computers \& Operations Research} 34(8),
2403--2435, 2007.

\bibitem{ropke2006} S.~Ropke, D.~Pisinger. An adaptive large neighborhood
search heuristic for the pickup and delivery problem with time windows.
\emph{Transportation Science} 40(4), 455--472, 2006.

\bibitem{savelsbergh1992} M.W.P.~Savelsbergh. The vehicle routing problem
with time windows: minimizing route duration.
\emph{ORSA Journal on Computing} 4(2), 146--154, 1992.

\bibitem{shaw1998} P.~Shaw. Using constraint programming and local search
methods to solve vehicle routing problems. \emph{Principles and Practice of
Constraint Programming} (CP 1998), LNCS 1520, 417--431, 1998.

\bibitem{vansteenwegen2009} P.~Vansteenwegen, W.~Souffriau,
G.~Vanden Berghe, D.~Van Oudheusden. Iterated local search for the team
orienteering problem with time windows.
\emph{Computers \& Operations Research} 36(12), 3281--3290, 2009.

\bibitem{vansteenwegen2011} P.~Vansteenwegen, W.~Souffriau,
D.~Van Oudheusden. The orienteering problem: a survey.
\emph{European Journal of Operational Research} 209(1), 1--10, 2011.

\bibitem{vidal2012} T.~Vidal, T.G.~Crainic, M.~Gendreau, N.~Lahrichi,
W.~Rei. A hybrid genetic algorithm for multidepot and periodic vehicle
routing problems. \emph{Operations Research} 60(3), 611--624, 2012.

\bibitem{vidal2022} T.~Vidal. Hybrid genetic search for the CVRP:
open-source implementation and SWAP* neighborhood.
\emph{Computers \& Operations Research} 140, 105643, 2022.

\end{thebibliography}

\end{document}
